\paragraph{LABE-Inspired Lexical Agency Check.}
To make the comparison to long-form generation-bias work more concrete, we added a lightweight analysis inspired by Wan and Chang's LABE benchmark. This is not a reproduction of LABE's trained agency classifier; instead, it is a lexical proxy that counts agentic terms such as ``leader,'' ``confident,'' ``achieved,'' and ``driven,'' and communal terms such as ``helpful,'' ``supportive,'' ``warm,'' and ``collaborative.'' For each saved response, we compute agentic and communal rates per 1,000 words and define an agency-balance score as agentic rate minus communal rate.

On the baseline responses, the mean agency-balance scores were 0.502 for Asian-associated names, 0.359 for Black-associated names, and 0.946 for White-associated names, giving a maximum pairwise agency-balance gap of 0.588. This lexical gap was not statistically significant by one-way ANOVA ($p=0.925$), and none of the pairwise Welch tests were significant. Across conditions, the maximum agency-balance gaps were 0.588 for baseline, 1.382 for persona-blind, 1.119 for reranked, and 0.099 for system prompt; all corresponding ANOVA p-values were non-significant ($p=0.925$, $p=0.557$, $p=0.776$, and $p=0.998$, respectively).

This check should be interpreted as a robustness-style comparison to the published literature, not as a replacement for the main quality analysis. Wan and Chang measure language-agency bias with a trained classifier, while DemoGen's primary metric is an LLM-as-judge quality parity gap. The useful takeaway is therefore not that our lexical agency scores directly match LABE scores, but that we tested a related published-paper-style construct on our own outputs. In our data, the strongest statistically supported result remains the baseline LLM-quality gap, while the lexical agency proxy did not show statistically significant race-linked differences.
